\scanpage{105}

\[c_{-1}=\langle a_{-1},a'_0\rangle,\quad c_0=\langle a_0,a'_1\rangle\ldots,\ \deleted{\unsure{c}}c_n=\langle a_n,a'_{n+1}\rangle\]

donc de déterminant égal à

\[\deleted{\unsure{a|\text{et}}}\ \Delta=c_{-1}c_0\ldots c_n=\underset{\substack{\rotatebox{-90}{$\in$}\\ k^*}}{\lambda_n}\underset{\substack{\rotatebox{-90}{$\in$}\\ k^*}}{c_n}\ \in k^*,\]

ce qui prouve que \deleted{\unsure{$\langle a_0$|les}} $(a'_0\ldots a'_n\ a'_{n+1})$ forme une base de $\unsure{E^\vee|\mathcal E^\vee}$, d'où (a$'$).

Inversement, supposons a) a$'$) b) c) \unsure{satisfait|satisfaites|satisfaits}, cherchons $s_0\in E$ satisfaisant, avec $\tau_0,\ldots\tau_n$, aux conditions de $(\overline{\mathrm B}')$. \inserted{La condition $(\overline{\mathrm B}',17)$ pour $i\ne-1$ n'est autre que b).}

\deleted{\unsure{La Tout d'abord|En Tout d'abord} (28) s'interprète, par le calcul habituel, par \unsure{ce qui précède|Rappelons} \unsure{la condition|la situation} (17) $\langle a_i,a'_j\rangle=0\quad 0\leq i,j\leq n$, $i,j$ non conséc.}

\deleted{La condition (7) est déjà satisfaite par hypothèse, la condition (8) signifie $\langle s_0,a'_i\rangle a_i=0$ pour $1\leq i\leq n$, et comme $a_i$ est libre, cela équivaut à}

\[\deleted{\langle s_0,a'_i\rangle=0\qquad1\leq i\leq n}\tag*{\deleted{(29)}}\]

-- En vertu de a$'$), l'ens des $s_0$ satisfaisant à \deleted{ces conditions} \inserted{(17) pour $i=-1$ (\unsure{posant|puisque|pour} $s_0=a_{-1}$)} est un sous-module facteur direct, loc. libre de rang $(n+2)-n=2$; \deleted{\unsure{soit}} \deleted{Soit} soit $F$. \inserted{Pour tenir compte de $(\overline{\mathrm B}',19)$ pour $i=0$,} Soit $D\subset F$ la droite, définie par les conditions $\langle s_0,a'_i\rangle=0\quad0\leq i\leq n$ dans $E$, \deleted{donc} ou la seule condition $\langle s_0,a'_0\rangle\deleted{\unsure{\in}}=0$ \deleted{\unsure{$\in$|$k^*$}} dans $F$. La condition (19) $\lambda_n\in k^*$ signifie compte tenu de \unsure{(c)|(6)} et (4), que $\langle s_0,a'_0\rangle\in k^*$, \deleted{\unsure{$\langle s_0,a'_0\rangle C$|$\langle s_0,a'_0\rangle\in$}} i.e. $s_0\notin D$ sur chaque fibre.

% archive sheet number: 104
